\chapter{前线轨道与反应活性}\label{ch:b2:frontier-orbitals}

放在架子上的一瓶环戊二烯会慢慢变成它的二聚体：
两个相同的二烯分子结合成一个桥式双环化合物，并且
只生成可能的立体异构体中的一种。本科一年级分册的弯箭头
无法预测是哪一种。两个轨道却可以：一个分子的最高占据轨道
与另一个分子的最低未占轨道。本章
把\cref{ch:b2:diatomic-mos}中的两能级相互作用
发展为研究反应活性的工具——亲核试剂进攻哪个原子、从哪一侧进攻，
以及二烯与烯烃为何能在一步反应中结合成环。

\begin{recall}
\Cref{ch:b2:diatomic-mos}：两个相互作用的轨道，四电子排斥。
\Cref{ch:b2:huckel}：休克尔能级与共轭体系的系数。
本科一年级分册：亲核试剂与亲电试剂、弯箭头、动力学
控制、立体专一反应、立体\omterm{def:b2:continuous-reactors:selectivity}{选择性}反应与区域\omterm{def:b2:continuous-reactors:selectivity}{选择性}反应，
烯烃的亲电加成、亲核取代。
\end{recall}

\section{两个轨道再探}

两个分子相互靠近时，它们的轨道发生重叠并相互作用。大部分
相互作用较弱，因为两\omterm{def:b2:atomic-orbitals:orbital-energy}{轨道能量}相差较大；
此时两能级结果的微扰形式便已足够。

\begin{theorem}[两能级的弱相互作用]\label{thm:b2:frontier-orbitals:perturbation}
能量分别为 $\varepsilon_1 < \varepsilon_2$ 的两个轨道，由\omterm{def:b2:diatomic-mos:integrals}{共振积分}
$\beta$ 耦合，且 $|\beta| \ll \Delta\varepsilon = \varepsilon_2 - \varepsilon_1$（忽略重叠），
则两能级移至
\[
  \varepsilon_1 - \frac{\beta^2}{\Delta\varepsilon}, \qquad \varepsilon_2 + \frac{\beta^2}{\Delta\varepsilon}
\]
在 $\beta^2/\Delta\varepsilon^2$ 的一阶近似下：下能级降低，上能级升高，
幅度相同，能隙越大，移动越小。
\end{theorem}

\begin{proof}
由\cref{thm:b2:diatomic-mos:heteronuclear}，精确能级为 $\bar\varepsilon \mp
\sqrt{\Delta\varepsilon^2/4 + \beta^2}$。令 $u = 4\beta^2/\Delta\varepsilon^2 \ll 1$，则 $\sqrt{\Delta\varepsilon^2/4
+ \beta^2} = \frac{\Delta\varepsilon}{2}\sqrt{1 + u} \approx \frac{\Delta\varepsilon}{2}(1 + u/2) =
\frac{\Delta\varepsilon}{2} + \frac{\beta^2}{\Delta\varepsilon}$，且 $\bar\varepsilon \mp \Delta\varepsilon/2 = \varepsilon_1,
\varepsilon_2$。
\end{proof}

\begin{proposition}[两电子相吸，四电子相斥]\label{prop:b2:frontier-orbitals:four-electron}
若两个相互作用的轨道共容纳两个电子，则体系稳定化，
能量降低约 $2\beta^2/\Delta\varepsilon$。若容纳四个电子，则体系去稳定化。
\end{proposition}

\begin{proof}
两个电子进入下能级，能量降低 $2\beta^2/\Delta\varepsilon$。四个电子时，上能级
也被填满；在\cref{thm:b2:frontier-orbitals:perturbation}的精度下，两项位移
相互抵消。计入重叠后，上能级的升高量大于下能级的降低量
（见\cref{prop:b2:diatomic-mos:asymmetry}），因此为净排斥。
\end{proof}

\begin{omfigure}
\begin{tikzpicture}[font=\small]
  \begin{scope}
    \pic (a) at (0,0) {omlevel={ud}{0.7}};
    \pic (b) at (3,1.4) {omlevel={empty}{0.7}};
    \pic (lo) at (1.5,-0.5) {omlevel={ud}{0.7}};
    \pic (hi) at (1.5,1.9) {omlevel={empty}{0.7}};
    \draw[omcorr, densely dashed] (a-r) -- (lo-l) (a-r) -- (hi-l) (b-l) -- (lo-r) (b-l) -- (hi-r);
    \node[below=6pt] at (1.5,-0.6) {两个电子：稳定化};
    \node[left] at (a-l) {已占据};
    \node[right] at (b-r) {空};
  \end{scope}
  \begin{scope}[xshift=7cm]
    \pic (a) at (0,0) {omlevel={ud}{0.7}};
    \pic (b) at (3,0.6) {omlevel={ud}{0.7}};
    \pic (lo) at (1.5,-0.5) {omlevel={ud}{0.7}};
    \pic (hi) at (1.5,1.25) {omlevel={ud}{0.7}};
    \draw[omcorr, densely dashed] (a-r) -- (lo-l) (a-r) -- (hi-l) (b-l) -- (lo-r) (b-l) -- (hi-r);
    \node[below=6pt] at (1.5,-0.6) {四个电子：排斥};
    \node[left] at (a-l) {已占据};
    \node[right] at (b-r) {已占据};
  \end{scope}
\end{tikzpicture}

{\small 相邻分子两个轨道的相互作用。已占据轨道
与空轨道相对时产生净稳定化；两个已占据轨道相互排斥，
上能级的升高量大于下能级的降低量。这种排斥是
分子间空间位阻的来源。}
\end{omfigure}

\section{前线轨道}

\begin{definition}[前线轨道]\label{def:b2:frontier-orbitals:frontier}
分子的\emph{HOMO}\index{HOMO}（最高占据\omterm{def:b2:diatomic-mos:lcao}{分子轨道}）和
\emph{LUMO}\index{LUMO}（最低未占\omterm{def:b2:diatomic-mos:lcao}{分子轨道}）称为其
\emph{前线轨道}\index{前线轨道}。
\end{definition}

\begin{proposition}[主导相互作用]\label{prop:b2:frontier-orbitals:dominant}
在两个闭壳层分子之间，起主要作用的稳定化相互作用
是一个分子的 \omterm{def:b2:frontier-orbitals:frontier}{HOMO} 与另一个分子的 \omterm{def:b2:frontier-orbitals:frontier}{LUMO} 之间的相互作用；在这两类配对中，
能隙较小者占主导。亲核试剂通过其 \omterm{def:b2:frontier-orbitals:frontier}{HOMO} 反应，
亲电试剂通过其 \omterm{def:b2:frontier-orbitals:frontier}{LUMO} 反应。
\end{proposition}

\begin{proof}[论证]
已占据–已占据配对相互排斥（见\cref{prop:b2:frontier-orbitals:four-electron}），而空–空配对
不含电子：只有已占据–空配对能稳定化，各降低能量 $2\beta^2/\Delta\varepsilon$。
\omterm{def:b2:frontier-orbitals:frontier}{HOMO} 是最高的已占能级，\omterm{def:b2:frontier-orbitals:frontier}{LUMO} 是最低的空能级，因此
\omterm{def:b2:frontier-orbitals:frontier}{HOMO}–\omterm{def:b2:frontier-orbitals:frontier}{LUMO} 配对在其中能隙最小；由\cref{thm:b2:frontier-orbitals:perturbation}可知，
能隙最小者给出最大的项。
\end{proof}

\begin{definition}[电荷控制与轨道控制]\label{def:b2:frontier-orbitals:control}
两反应物所带电荷之间的吸引决定反应位置时，反应
受\emph{电荷控制}\index{电荷控制}；
\omterm{def:b2:frontier-orbitals:frontier}{前线轨道}重叠决定反应位置时，反应受\emph{轨道控制}\index{轨道控制}，
进攻发生在系数最大的位置。
\end{definition}

\begin{definition}[硬物种与软物种]\label{def:b2:frontier-orbitals:hard-soft}
\emph{硬亲核试剂}\index{硬亲核试剂}或\emph{硬亲电试剂}
\index{硬亲电试剂}体积小，电荷集中，具有 \omterm{def:b2:frontier-orbitals:frontier}{HOMO}（亲电试剂则为 \omterm{def:b2:frontier-orbitals:frontier}{LUMO}）
远离反应伙伴的相应轨道，受\omterm{def:b2:frontier-orbitals:control}{电荷控制}。
\emph{软亲核试剂}\index{软亲核试剂}或\emph{软亲电试剂}
\index{软亲电试剂}体积大，易极化，\omterm{def:b2:frontier-orbitals:frontier}{HOMO} 高
（亲电试剂则为低 \omterm{def:b2:frontier-orbitals:frontier}{LUMO}），受\omterm{def:b2:frontier-orbitals:control}{轨道控制}。硬者优先与硬者反应，
软者与软者反应。
\end{definition}

氢氧根离子和氟离子是\omterm{def:b2:frontier-orbitals:hard-soft}{硬亲核试剂}，碘离子和硫醇盐
是\omterm{def:b2:frontier-orbitals:hard-soft}{软亲核试剂}；质子和碳正离子是\omterm{def:b2:frontier-orbitals:hard-soft}{硬亲电试剂}，卤代烷中的碳
则是\omterm{def:b2:frontier-orbitals:hard-soft}{软亲电试剂}。

\begin{definition}[两可亲核试剂]\label{def:b2:frontier-orbitals:ambident}
\emph{两可亲核试剂}\index{两可亲核试剂}含有两个亲核原子，
这两个原子通过共轭相连，各自都可以进攻亲电试剂。
\end{definition}

氰根离子被\omterm{def:b2:frontier-orbitals:hard-soft}{硬亲电试剂}进攻时，发生在负电荷较大的氮原子上，
被\omterm{def:b2:frontier-orbitals:hard-soft}{软亲电试剂}进攻时，发生在 \omterm{def:b2:frontier-orbitals:frontier}{HOMO} 系数较大的碳原子上：
卤代烷生成腈 \ce{R-C#N}。烯醇负离子已在
\cref{ch:b2:enolates-aldol}中见过，它们同样在碳或氧上发生反应。

\begin{method}[前线轨道分析]\label{met:b2:frontier-orbitals:fmo}
\begin{enumerate}
  \item 确定亲核试剂（\omterm{def:b2:frontier-orbitals:frontier}{HOMO} 高）和亲电试剂（\omterm{def:b2:frontier-orbitals:frontier}{LUMO} 低）。
  \item 画出两个\omterm{def:b2:frontier-orbitals:frontier}{前线轨道}，标明其系数与符号。
  \item 反应发生在重叠最大处：系数最大的原子，
        符号相同的瓣相互面对。
  \item 进攻方向由 \omterm{def:b2:frontier-orbitals:frontier}{LUMO} 的形状决定（垂直于
        \ce{C=O} 平面，在 \ce{C-X} 键背面）。
  \item HOMO--LUMO 能隙越小，反应越快（在其他条件相同时）。
\end{enumerate}
\end{method}

进攻方向由 \omterm{def:b2:frontier-orbitals:frontier}{LUMO} 决定。羰基的 \omterm{def:b2:frontier-orbitals:frontier}{LUMO} 是其
$\pi^*$ 轨道，碳上的瓣较大，且垂直于该平面：亲核试剂从平面上方或下方接近
碳原子，而不沿 \ce{C=O} 轴方向。卤代烷的 \omterm{def:b2:frontier-orbitals:frontier}{LUMO} 是
$\sigma^*$ 轨道，属于 \ce{C-X} 键，其在碳上的大瓣背向 X：
亲核试剂从背面进攻，碳原子发生构型翻转
——即本科一年级分册中双分子取代的立体化学。

\section{第尔斯–阿尔德反应}

\begin{definition}[协同反应]\label{def:b2:frontier-orbitals:concerted}
\emph{协同反应}\index{协同反应}是指在一步之内，
经一个过渡态完成全部成键与断键，不经过中间体。
\end{definition}

\begin{definition}[第尔斯–阿尔德反应]\label{def:b2:frontier-orbitals:diels-alder}
\emph{第尔斯–阿尔德反应}\index{第尔斯–阿尔德反应}是共轭的
\emph{双烯体}\index{双烯体}在 s-顺式构象下与烯烃或炔烃
（即\emph{亲双烯体}\index{亲双烯体}）的协同加成，形成六元环，
并生成两个新的 $\sigma$ 键和一个新的 $\pi$ 键。
\end{definition}

\begin{omfigure}
\begin{tikzpicture}[font=\small, >={Stealth}, scale=1.25]
  % diene s-cis: C1 (0,1.6) C2 (0.6,2.4) C3 (1.6,2.4) C4 (2.2,1.6); dienophile below
  \draw[thick] (0,1.6) -- (0.6,2.4) -- (1.6,2.4) -- (2.2,1.6);
  \draw[thick] (0.2,1.62) -- (0.65,2.22);
  \draw[thick] (2.0,1.62) -- (1.55,2.22);
  \draw[thick] (0.2,0) -- (2.0,0); \draw[thick] (0.35,0.14) -- (1.85,0.14);
  \node[font=\scriptsize, anchor=east] at (-0.05,1.6) {1};
  \node[font=\scriptsize, anchor=south east] at (0.5,2.42) {2};
  \node[font=\scriptsize, anchor=south] at (1.6,2.42) {3};
  \node[font=\scriptsize, anchor=west] at (2.25,1.6) {4};
  % arrows: dienophile pi -> C1..C(dienophile); C1=C2 pi -> C2-C3; C3=C4 pi -> C4-C
  \draw[->, thick, omProp] (1.1,0.22) .. controls (0.9,0.9) and (0.3,0.9) .. (0.12,1.42);
  \draw[->, thick, omProp] (0.45,1.95) .. controls (0.6,2.75) and (1.0,2.8) .. (1.1,2.48);
  \draw[->, thick, omProp] (1.78,1.92) .. controls (2.5,1.6) and (2.4,0.6) .. (2.05,0.2);
  \node at (3.3,1.2) {$\longrightarrow$};
  \begin{scope}[xshift=4.2cm]
    \draw[thick] (0,1.6) -- (0.6,2.4) -- (1.6,2.4) -- (2.2,1.6) -- (2.0,0) -- (0.2,0) -- cycle;
    \draw[thick] (0.72,2.26) -- (1.48,2.26);
  \end{scope}
\end{tikzpicture}

{\small 丁二烯与乙烯的\omterm{def:b2:frontier-orbitals:diels-alder}{第尔斯–阿尔德反应}，用弯箭头表示：
三对电子同时移动，经一个环状协同步骤完成；两个 $\sigma$ 键
在\omterm{def:b2:frontier-orbitals:diels-alder}{双烯体}两端形成，$\pi$ 键移至其中心。}
\end{omfigure}

主导的相互作用是\omterm{def:b2:frontier-orbitals:diels-alder}{双烯体}的 \omterm{def:b2:frontier-orbitals:frontier}{HOMO}（$\psi_2$）与
\omterm{def:b2:frontier-orbitals:diels-alder}{亲双烯体}的 \omterm{def:b2:frontier-orbitals:frontier}{LUMO}（$\pi^*$）。\omterm{def:b2:frontier-orbitals:diels-alder}{双烯体}两端的瓣与\omterm{def:b2:frontier-orbitals:diels-alder}{亲双烯体}的两个碳上的瓣
在\omterm{def:b2:frontier-orbitals:diels-alder}{亲双烯体}靠近\omterm{def:b2:frontier-orbitals:diels-alder}{双烯体}的一个面时，符号相匹配：
两个新键可同时在两反应物各自的同一面形成。

\begin{omfigure}
\begin{tikzpicture}[font=\small]
  \foreach \x/\c in {0/0.602, 1/0.372, 2/-0.372, 3/-0.602} {
    \pic at (\x,2) {omp={1.5*\c}{90}};
  }
  \draw[black!45] (-0.2,2) -- (3.2,2);
  \node[anchor=west] at (3.6,2) {双烯体的 HOMO（$\psi_2$）};
  \pic at (0,0) {omp={-0.85}{90}}; \pic at (3,0) {omp={0.85}{90}};
  \draw[black!45] (-0.2,0) -- (3.2,0);
  \node[anchor=west] at (3.6,0) {亲双烯体的 LUMO（$\pi^*$）};
  \draw[omProp, thick, densely dashed] (0,1.45) -- (0,0.6);
  \draw[omProp, thick, densely dashed] (3,1.45) -- (3,0.6);
  \node[font=\scriptsize, omProp, anchor=east] at (-0.15,1.0) {同相};
  \node[font=\scriptsize, omProp, anchor=west] at (3.15,1.0) {同相};
\end{tikzpicture}

{\small 丁二烯与乙烯的\omterm{def:b2:frontier-orbitals:frontier}{前线轨道}，瓣的大小与其
休克尔系数成正比。两者相对时，\omterm{def:b2:frontier-orbitals:diels-alder}{双烯体} \omterm{def:b2:frontier-orbitals:frontier}{HOMO} 的端部瓣与
\omterm{def:b2:frontier-orbitals:diels-alder}{亲双烯体} \omterm{def:b2:frontier-orbitals:frontier}{LUMO} 的瓣在两端符号相同：两个
$\sigma$ 键可在两反应物各自的同一面同时形成。}
\end{omfigure}

\begin{proposition}[立体专一性]\label{prop:b2:frontier-orbitals:da-stereospecific}
\omterm{def:b2:frontier-orbitals:diels-alder}{第尔斯–阿尔德反应}具有立体专一性：\omterm{def:b2:frontier-orbitals:diels-alder}{亲双烯体}上的顺式取代基
在产物中仍为顺式，反式取代基仍保持反式。
\end{proposition}

\begin{proof}
该反应为\omterm{def:b2:frontier-orbitals:concerted}{协同反应}，两个新键都在\omterm{def:b2:frontier-orbitals:diels-alder}{亲双烯体}的
同一面形成：其两个碳原子彼此从不发生相对旋转，两者上取代基的相对位置
被带入环中。
\end{proof}

\begin{proposition}[区域选择性]\label{prop:b2:frontier-orbitals:da-regio}
对于不对称的反应物，主要产物是将\omterm{def:b2:frontier-orbitals:diels-alder}{双烯体}中
\omterm{def:b2:frontier-orbitals:frontier}{HOMO} 中系数最大的原子与\omterm{def:b2:frontier-orbitals:diels-alder}{亲双烯体}中
\omterm{def:b2:frontier-orbitals:frontier}{LUMO} 中系数最大的原子连接而成。
\end{proposition}

\begin{proof}[论证]
在过渡态中，两个新键的形成程度尚不相同；
当较大系数的原子相互重叠时，稳定化能 $2\beta^2/\Delta\varepsilon$ 更大（$\beta$
随重叠原子系数之积增大），因此
该取向在较低能量下即可达到。
\end{proof}

\begin{omfigure}
\begin{tikzpicture}
\begin{axis}[ybar, width=0.47\linewidth, height=4.6cm, ymin=-0.8, ymax=0.8,
  xtick={0,1,2,3,4}, xticklabels={O,C1,C2,C3,C4}, axis x line=bottom, axis y line=left,
  extra y ticks={0}, extra y tick style={grid=major, grid style={black!50}}, enlarge x limits=0.15,
  tick label style={font=\scriptsize}, ylabel={系数}, label style={font=\small},
  title={\small 1-甲氧基丁二烯的 HOMO}, bar width=8pt]
\addplot[fill=omDef!50, draw=omDef] table[x=atom, y=c] {figdata/out/bachelor-2-huckel-c-diene.dat};
\end{axis}
\end{tikzpicture}\hfill
\begin{tikzpicture}
\begin{axis}[ybar, width=0.47\linewidth, height=4.6cm, ymin=-0.8, ymax=0.8,
  xtick={1,2,3,4}, xticklabels={C$\beta$,C$\alpha$,C(O),O}, axis x line=bottom, axis y line=left,
  extra y ticks={0}, extra y tick style={grid=major, grid style={black!50}}, enlarge x limits=0.15,
  tick label style={font=\scriptsize}, ylabel={系数}, label style={font=\small},
  title={\small 丙烯醛的 LUMO}, bar width=8pt]
\addplot[fill=omThm!45, draw=omThm] table[x=atom, y=c] {figdata/out/bachelor-2-huckel-c-dienophile.dat};
\end{axis}
\end{tikzpicture}

{\small 休克尔模型，含杂原子参数（醚氧：$\alpha_{\ce{O}} = \alpha + 2\beta$，
$\beta_{\ce{CO}} = 0.8\beta$；羰基氧：$\alpha_{\ce{O}} = \alpha + \beta$，$\beta_{\ce{CO}} =
\beta$；此为模型）。\omterm{def:b2:frontier-orbitals:diels-alder}{双烯体}的 \omterm{def:b2:frontier-orbitals:frontier}{HOMO} 在 C4 上最大，即
远离甲氧基的一端（0.60 对 0.51）；\omterm{def:b2:frontier-orbitals:diels-alder}{亲双烯体}的 \omterm{def:b2:frontier-orbitals:frontier}{LUMO} 在
$\beta$ 碳上最大（0.66）。两者成键后，甲氧基与醛基
位于环上相邻的碳原子上。}
\end{omfigure}

给电子基团使\omterm{def:b2:frontier-orbitals:diels-alder}{双烯体}的 \omterm{def:b2:frontier-orbitals:frontier}{HOMO} 升高（$m = 0.48$，原为 0.62），吸电子基团
使\omterm{def:b2:frontier-orbitals:diels-alder}{亲双烯体}的 \omterm{def:b2:frontier-orbitals:frontier}{LUMO} 降低（$m = -0.35$，原为 $-1$）：能隙变窄，因此这类
反应物对的反应速率远快于丁二烯与乙烯的反应。富电子的\omterm{def:b2:frontier-orbitals:diels-alder}{双烯体}与
缺电子的\omterm{def:b2:frontier-orbitals:diels-alder}{亲双烯体}是典型的组合。

\begin{definition}[内型与外型加合物]\label{def:b2:frontier-orbitals:endo}
环状\omterm{def:b2:frontier-orbitals:diels-alder}{双烯体}与带有不饱和取代基的\omterm{def:b2:frontier-orbitals:diels-alder}{亲双烯体}加成时，
\emph{内型加合物}\index{内型加合物}中该取代基指向新形成的
双键（过渡态中位于\omterm{def:b2:frontier-orbitals:diels-alder}{双烯体}之下），\emph{外型加合物}\index{外型加合物}中该取代基则
背离该双键。\emph{内型规则}\index{内型规则}指出：
内型加合物生成得更快。
\end{definition}

\begin{proposition}[内型规则]\label{prop:b2:frontier-orbitals:endo}
在动力学控制下，\omterm{def:b2:frontier-orbitals:endo}{内型加合物}为主要产物，尽管\omterm{def:b2:frontier-orbitals:endo}{外型加合物}
往往更稳定。
\end{proposition}

\admitted

通常的解释是：在内型接近中，取代基 $\pi$ 体系的瓣与
\omterm{def:b2:frontier-orbitals:diels-alder}{双烯体} \omterm{def:b2:frontier-orbitals:frontier}{HOMO} 中心原子的瓣之间存在次级重叠，
它降低了内型过渡态的能量，但并不形成键。该规则是
经验性的，存在例外；本科三年级分册将借助
轨道的对称性来处理这些反应。

\begin{omfigure}
\begin{tikzpicture}[font=\small]
  \foreach \xs/\mode in {0/内型, 5/外型} {
    \begin{scope}[xshift=\xs cm]
      % cyclopentadiene seen from the side: C1, C2=C3, C4, CH2 bridge on top
      \draw[thick] (0,2) -- (0.6,2.5) -- (1.8,2.5) -- (2.4,2);
      \draw[thick] (0.75,2.38) -- (1.65,2.38);
      \draw[thick] (0,2) -- (1.2,3.2) -- (2.4,2);
      \node[font=\scriptsize, anchor=south] at (1.2,3.2) {\ce{CH2}};
      % dienophile
      \draw[thick] (0.2,0) -- (2.2,0); \draw[thick] (0.35,0.1) -- (2.05,0.1);
      \draw[omProp, densely dashed, thick] (0,1.9) -- (0.2,0.15);
      \draw[omProp, densely dashed, thick] (2.4,1.9) -- (2.2,0.15);
      \node[anchor=north] at (1.2,-1.3) {\mode};
    \end{scope}
  }
  % endo: carbonyls tucked under the diene
  \draw[thick] (0.2,0) -- (0.75,0.95) (2.2,0) -- (1.65,0.95);
  \node[font=\scriptsize, fill=white, inner sep=1pt] at (0.75,1.1) {C=O};
  \node[font=\scriptsize, fill=white, inner sep=1pt] at (1.65,1.1) {C=O};
  \draw[omThm, densely dotted, thick] (0.75,1.25) -- (0.65,2.35);
  \draw[omThm, densely dotted, thick] (1.65,1.25) -- (1.75,2.35);
  % exo: carbonyls pointing away
  \draw[thick] (5.2,0) -- (4.95,-0.85) (7.2,0) -- (7.45,-0.85);
  \node[font=\scriptsize] at (4.95,-1.0) {C=O};
  \node[font=\scriptsize] at (7.45,-1.0) {C=O};
\end{tikzpicture}

{\small 顺丁烯二酸酐与环戊二烯的内型与外型接近方式，示意。
形成中的键用虚线表示。在内型接近中，羰基位于
\omterm{def:b2:frontier-orbitals:diels-alder}{双烯体}之下，其 $\pi$ 轨道与\omterm{def:b2:frontier-orbitals:diels-alder}{双烯体}的中心碳原子重叠（点线，次级
重叠）；在外型接近中，羰基则背离\omterm{def:b2:frontier-orbitals:diels-alder}{双烯体}。}
\end{omfigure}

\begin{method}[画出第尔斯–阿尔德反应的产物]\label{met:b2:frontier-orbitals:diels-alder}
\begin{enumerate}
  \item 使\omterm{def:b2:frontier-orbitals:diels-alder}{双烯体}处于 s-顺式构象，\omterm{def:b2:frontier-orbitals:diels-alder}{亲双烯体}正对其两端。
  \item 画出从\omterm{def:b2:frontier-orbitals:diels-alder}{双烯体}两端到\omterm{def:b2:frontier-orbitals:diels-alder}{亲双烯体}碳原子的两个新 $\sigma$ 键，
        以及\omterm{def:b2:frontier-orbitals:diels-alder}{双烯体} C2 与 C3 之间的新 $\pi$ 键。
  \item 保持\omterm{def:b2:frontier-orbitals:diels-alder}{亲双烯体}上的取代基位于同侧（顺式仍为顺式）。
  \item 对于不对称的反应物，应将轨道系数最大的原子相连（生成“邻位”或“对位”
        产物，分别对应 1-或 2-取代\omterm{def:b2:frontier-orbitals:diels-alder}{双烯体}）。
  \item 对于环状\omterm{def:b2:frontier-orbitals:diels-alder}{双烯体}，应使不饱和取代基处于内型位置。
\end{enumerate}
\end{method}

\begin{inthelab}[裂解二环戊二烯]
环戊二烯以二聚体形式出售，使用前需即时再生：将二聚体
加热至沸点（约 \qty{170}{\degreeCelsius}），于配有
\omterm{def:b2:liquid-vapour-diagrams:fractional-distillation}{分馏}柱的烧瓶中进行，二聚体在此分解为单体（即逆\omterm{def:b2:frontier-orbitals:diels-alder}{第尔斯–阿尔德
反应}），单体在 \qty{41}{\degreeCelsius} 馏出，收集于
冰浴冷却的接收器中。单体于室温下数小时内即再次二聚，应立即使用。
两种化合物均易燃且有害，操作须在通风橱中进行。
\end{inthelab}

\begin{safety}{\ghs[5mm]{GHS02}\,\ghs[5mm]{GHS06}\,\ghs[5mm]{GHS08}}
环戊二烯及其二聚体极易燃，吸入有毒；顺丁烯二酸
酐具有腐蚀性，是呼吸道致敏物。它们须在通风橱中操作，
并戴手套和护目镜，远离火源。
% ledger: ghs:cyclopentadiene, ghs:dicyclopentadiene, ghs:maleic-anhydride, bp:cyclopentadiene, bp:dicyclopentadiene
\end{safety}

\section{练习}

\begin{exercise}[$\star$]\label{exo:b2:frontier-orbitals:1}
求乙烯、烯丙基负离子以及丁二烯的 \omterm{def:b2:frontier-orbitals:frontier}{HOMO} 与 \omterm{def:b2:frontier-orbitals:frontier}{LUMO}
（能量以休克尔单位表示）。
\end{exercise}

\begin{exercise}[$\star$]\label{exo:b2:frontier-orbitals:2}
判断下列物种属硬还是软：\ce{OH-}、\ce{I-}、\ce{H+}、硫醇盐 \ce{RS-}、碘甲烷中的
碳原子、\ce{F-}。
\end{exercise}

\begin{exercise}[$\star$]\label{exo:b2:frontier-orbitals:3}
下列哪些二烯能参与\omterm{def:b2:frontier-orbitals:diels-alder}{第尔斯–阿尔德反应}：1,3-丁二烯、
环戊二烯、1,4-戊二烯、（2\textit{Z},4\textit{Z}）-己-2,4-二烯
（以 s-顺式构象存在）？
\end{exercise}

\begin{exercise}[$\star$]\label{exo:b2:frontier-orbitals:4}
画出丁二烯与丙烯腈 \ce{CH2=CH-C#N} 的产物。
\end{exercise}

\begin{exercise}[$\star\star$]\label{exo:b2:frontier-orbitals:5}
两个轨道的能量分别为 $-10$ 和 \qty{-2}{eV}，与 $\beta = \qty{-1.0}{eV}$ 相互作用（练习数据）。
先计算两个电子的微扰稳定化能，再计算精确的稳定化能。
\end{exercise}

\begin{exercise}[$\star\star$]\label{exo:b2:frontier-orbitals:6}
氰根离子与碘甲烷反应生成乙腈 \ce{CH3-C#N}，而与一个
碳正离子在强电离性溶剂中反应时，则生成部分异腈 \ce{R-N#C}。请解释。
\end{exercise}

\begin{exercise}[$\star\star$]\label{exo:b2:frontier-orbitals:7}
利用本章的轨道系数，预测 1-甲氧基丁二烯
与丙烯醛反应的主要产物。
\end{exercise}

\begin{exercise}[$\star\star$]\label{exo:b2:frontier-orbitals:8}
画出环戊二烯与丙烯酸甲酯反应生成的\omterm{def:b2:frontier-orbitals:endo}{内型加合物}。
\end{exercise}

\begin{exercise}[$\star\star$]\label{exo:b2:frontier-orbitals:9}
为何顺丁烯二酸酐在室温下即可与环戊二烯反应，而乙烯
则需在加热加压条件下反应？
\end{exercise}

\begin{exercise}[$\star\star\star$]\label{exo:b2:frontier-orbitals:10}
环戊二烯分别与马来酸二甲酯（顺式二酯）和富马酸二甲酯
（反式二酯）反应。画出两种产物，并说明各自给出几种立体异构体。
\end{exercise}

\begin{exercise}[$\star\star\star$]\label{exo:b2:frontier-orbitals:11}
\omterm{def:b2:frontier-orbitals:diels-alder}{第尔斯–阿尔德反应}满足 $\Delta_r H^\circ < 0$ 和 $\Delta_r S^\circ < 0$。请解释这些正负号
的含义，并说明为何逆反应在高温下变得有利。
\end{exercise}

\begin{exercise}[$\star\star\star$]\label{exo:b2:frontier-orbitals:12}
在肼 \ce{H2N-NH2} 中，每个氮原子均带一对孤电子对。利用四电子
相互作用，解释分子为何避免两对孤电子对
平行的构象。
\end{exercise}

\section{问题：裂解环戊二烯}

\begin{problem}[{周末问题——环戊二烯二聚化作为第尔斯–阿尔德
反应、裂解的热力学、环戊二烯与顺丁烯二酸酐的前线轨道，
以及它们加合物的立体化学}]
\label{pb:b2:frontier-orbitals:1}
数据：沸点，环戊二烯 \qty{41}{\degreeCelsius}，二环戊二烯约
\qty{170}{\degreeCelsius}。休克尔模型能量（仅为一种模型）：环戊二烯视为一个丁二烯单元，
其 \omterm{def:b2:frontier-orbitals:frontier}{HOMO} 为 $\alpha + 0.62\beta$，\omterm{def:b2:frontier-orbitals:frontier}{LUMO} 为 $\alpha - 0.62\beta$；顺丁烯二酸酐的 \omterm{def:b2:frontier-orbitals:frontier}{LUMO}
约为 $\alpha - 0.35\beta$。
% ledger: bp:cyclopentadiene, bp:dicyclopentadiene

\medskip
\noindent\textbf{第一部分——二聚体。}
\begin{enumerate}
  \item 画出环戊二烯，并说明其双烯单元被固定为 s-顺式构象。
  \item 二聚化过程中，一个分子作为\omterm{def:b2:frontier-orbitals:diels-alder}{双烯体}，另一个分子是什么？
  \item 画出二聚体，并对新形成的键编号。
  \item 室温下生成哪一种加合物（内型或外型），为什么？
  \item 该二聚化是否具有立体专一性？请结合反应机理说明。
  \item 为何该二聚化无需催化剂？
\end{enumerate}

\medskip
\noindent\textbf{第二部分——裂解的热力学。}
\begin{enumerate}[resume]
  \item 由所形成和断裂的化学键，判断二聚化的 $\Delta_r H^\circ$ 的符号，
        并说明依据。
  \item 判断 $\Delta_r S^\circ$ 的符号。
  \item 写出 $\Delta_r G^\circ(T)$，并证明它在温度 $T_i$ 处改变符号。
  \item 为何二聚体在室温下稳定，而单体在加热时更有利？
  \item 在裂解装置中，为何将单体一边生成一边蒸馏出来，能推动
        反应进行？
  \item 为何单体必须保存在低温下并尽快使用？
\end{enumerate}

\medskip
\noindent\textbf{第三部分——\omterm{def:b2:frontier-orbitals:frontier}{前线轨道}。}
\begin{enumerate}[resume]
  \item 考虑环戊二烯与自身反应，以
        $|\beta|$ 为单位，计算 HOMO--LUMO 能隙。
  \item 考虑环戊二烯与顺丁烯二酸酐反应，计算
        \omterm{def:b2:frontier-orbitals:diels-alder}{双烯体}的 \omterm{def:b2:frontier-orbitals:frontier}{HOMO} 与酸酐的 \omterm{def:b2:frontier-orbitals:frontier}{LUMO} 之间的能隙。
  \item 哪一对反应更快？请说明理由。
  \item 为何顺丁烯二酸酐中的两个羰基会降低其 \omterm{def:b2:frontier-orbitals:frontier}{LUMO}？
  \item 在内型接近方式中，顺丁烯二酸酐的哪一轨道
        提供次级重叠？
\end{enumerate}

\medskip
\noindent\textbf{第四部分——与顺丁烯二酸酐的加合物。}
\begin{enumerate}[resume]
  \item 画出\omterm{def:b2:frontier-orbitals:endo}{内型加合物}。
  \item 标出其中的手性碳原子。
  \item 该加合物是否具有手性？（提示：寻找对称面。）
  \item 为何原\omterm{def:b2:frontier-orbitals:diels-alder}{亲双烯体}上的两个氢原子
        会位于环的同一侧？
  \item \omterm{def:b2:frontier-orbitals:endo}{外型加合物}会是什么样子？
  \item 内型与\omterm{def:b2:frontier-orbitals:endo}{外型加合物}能否在室温下相互转化？为什么？
  \item 酸酐与水如何反应？会得到怎样的生成物？
  \item 写出环戊二烯与顺丁烯二酸酐的\omterm{def:b2:frontier-orbitals:endo}{内型加合物}中
        手性碳原子的数目。
\end{enumerate}
\end{problem}
